\documentclass[12pt,a4paper]{article}
\usepackage[utf8]{inputenc}
\usepackage[spanish]{babel}
\usepackage{geometry}
\usepackage{booktabs}
\usepackage{tabularx}
\usepackage{float}

\geometry{margin=2.5cm}

\title{Requerimientos Funcionales Seleccionados}
\author{}
\date{}

\begin{document}

\maketitle

\section{Requerimientos Funcionales Seleccionados}

La siguiente tabla presenta los requerimientos funcionales RF06, RF07, RF13, RF16 y RF39, organizados con el formato estándar utilizado en el proyecto.

\begin{table}[H]
\centering
\renewcommand{\arraystretch}{1.2}
\begin{tabularx}{\textwidth}{@{} l l X @{}}
\toprule
\textbf{Código} & \textbf{Nombre} & \textbf{Descripción} \\
\midrule

RF06 & Registro de ingresos &
El usuario podrá registrar manualmente sus ingresos indicando el monto, la fecha y una descripción opcional, tanto desde la aplicación móvil como desde la plataforma web. \\

RF07 & Registro de egresos &
El usuario podrá registrar egresos de forma manual o automática, especificando monto, fecha, categoría y descripción. \\

RF13 & Generación de reportes en PDF &
La plataforma web y la aplicación móvil deberán permitir la generación de reportes financieros en formato PDF, con opción de descarga o visualización en pantalla. \\

RF16 & Captura de foto del ticket &
La aplicación móvil debe permitir al usuario tomar una fotografía del ticket de compra utilizando la cámara del dispositivo para su posterior procesamiento. \\

RF39 & Carga de ticket digital &
El sistema debe permitir la carga de tickets en formato digital (PDF o imagen) desde el explorador web, procesándolos automáticamente con el módulo OCR. \\

\bottomrule
\end{tabularx}
\caption{Requerimientos Funcionales Seleccionados (RF06, RF07, RF13, RF16, RF39)}
\end{table}

\end{document}